% figs.tex —— 由 ch7/make_figs.py 產生，請勿手動修改。
% 重新產生：python ch7/make_figs.py（第 7 章，動畫 W701–W706）

\newcommand{\SevenFigGrid}{%
\begin{tikzpicture}[x=0.5cm,y=0.5cm,font=\scriptsize,baseline=(current bounding box.center)]
\draw[gray!30] (-2.5,-2.5) grid (2.5,2.5);
\draw[gray!60,->] (-2.7,0) -- (2.8,0); \draw[gray!60,->] (0,-2.7) -- (0,2.8);
\node[circle,draw=cBar,fill=cBar!10,inner sep=1.2pt,minimum size=11pt] at (-2,0) {$1$};
\node[circle,draw=cBar,fill=cBar!25,inner sep=1.2pt,minimum size=11pt] at (-1,-1) {$2$};
\node[circle,draw=cBar,fill=cBar!25,inner sep=1.2pt,minimum size=11pt] at (-1,1) {$2$};
\node[circle,draw=cBar,fill=cBar!10,inner sep=1.2pt,minimum size=11pt] at (0,-2) {$1$};
\node[circle,draw=cBar,fill=cBar!45,inner sep=1.2pt,minimum size=11pt] at (0,0) {$4$};
\node[circle,draw=cBar,fill=cBar!10,inner sep=1.2pt,minimum size=11pt] at (0,2) {$1$};
\node[circle,draw=cBar,fill=cBar!25,inner sep=1.2pt,minimum size=11pt] at (1,-1) {$2$};
\node[circle,draw=cBar,fill=cBar!25,inner sep=1.2pt,minimum size=11pt] at (1,1) {$2$};
\node[circle,draw=cBar,fill=cBar!10,inner sep=1.2pt,minimum size=11pt] at (2,0) {$1$};
\node at (0,-3.2) {$\sum_{\vec r}C(2;\vec r)^2=36$};
\end{tikzpicture}}

\newcommand{\SevenFigGauss}{%
\begin{tikzpicture}[x=0.36cm,y=6.2cm,>=Stealth,font=\scriptsize,baseline=(current bounding box.center)]
\draw[->] (-6.4,0) -- (6.6,0) node[right] {$x$};
\draw[->] (0,0) -- (0,0.47);
\foreach \x in {-6,-3,3,6} {\draw (\x,0) -- (\x,-0.012); \node[below] at (\x,-0.01) {$\x$};}
\foreach \y in {0.1,0.2,0.3,0.4} {\draw (0,\y) -- (-0.2,\y); \node[left] at (-0.2,\y) {$\y$};}
\draw[cBar,densely dashed,thick] plot coordinates {(-6.000,0.0000) (-5.833,0.0000) (-5.667,0.0000) (-5.500,0.0000) (-5.333,0.0000) (-5.167,0.0000) (-5.000,0.0000) (-4.833,0.0000) (-4.667,0.0000) (-4.500,0.0000) (-4.333,0.0000) (-4.167,0.0001) (-4.000,0.0001) (-3.833,0.0003) (-3.667,0.0005) (-3.500,0.0009) (-3.333,0.0015) (-3.167,0.0027) (-3.000,0.0044) (-2.833,0.0072) (-2.667,0.0114) (-2.500,0.0175) (-2.333,0.0262) (-2.167,0.0382) (-2.000,0.0540) (-1.833,0.0743) (-1.667,0.0995) (-1.500,0.1295) (-1.333,0.1640) (-1.167,0.2020) (-1.000,0.2420) (-0.833,0.2819) (-0.667,0.3194) (-0.500,0.3521) (-0.333,0.3774) (-0.167,0.3934) (0.000,0.3989) (0.167,0.3934) (0.333,0.3774) (0.500,0.3521) (0.667,0.3194) (0.833,0.2819) (1.000,0.2420) (1.167,0.2020) (1.333,0.1640) (1.500,0.1295) (1.667,0.0995) (1.833,0.0743) (2.000,0.0540) (2.167,0.0382) (2.333,0.0262) (2.500,0.0175) (2.667,0.0114) (2.833,0.0072) (3.000,0.0044) (3.167,0.0027) (3.333,0.0015) (3.500,0.0009) (3.667,0.0005) (3.833,0.0003) (4.000,0.0001) (4.167,0.0001) (4.333,0.0000) (4.500,0.0000) (4.667,0.0000) (4.833,0.0000) (5.000,0.0000) (5.167,0.0000) (5.333,0.0000) (5.500,0.0000) (5.667,0.0000) (5.833,0.0000) (6.000,0.0000)};
\draw[cWalk,densely dashed,thick] plot coordinates {(-6.000,0.0000) (-5.833,0.0001) (-5.667,0.0001) (-5.500,0.0001) (-5.333,0.0002) (-5.167,0.0004) (-5.000,0.0005) (-4.833,0.0008) (-4.667,0.0012) (-4.500,0.0018) (-4.333,0.0026) (-4.167,0.0037) (-4.000,0.0052) (-3.833,0.0072) (-3.667,0.0098) (-3.500,0.0132) (-3.333,0.0175) (-3.167,0.0230) (-3.000,0.0297) (-2.833,0.0379) (-2.667,0.0477) (-2.500,0.0591) (-2.333,0.0723) (-2.167,0.0872) (-2.000,0.1038) (-1.833,0.1217) (-1.667,0.1409) (-1.500,0.1607) (-1.333,0.1809) (-1.167,0.2007) (-1.000,0.2197) (-0.833,0.2371) (-0.667,0.2524) (-0.500,0.2650) (-0.333,0.2744) (-0.167,0.2801) (0.000,0.2821) (0.167,0.2801) (0.333,0.2744) (0.500,0.2650) (0.667,0.2524) (0.833,0.2371) (1.000,0.2197) (1.167,0.2007) (1.333,0.1809) (1.500,0.1607) (1.667,0.1409) (1.833,0.1217) (2.000,0.1038) (2.167,0.0872) (2.333,0.0723) (2.500,0.0591) (2.667,0.0477) (2.833,0.0379) (3.000,0.0297) (3.167,0.0230) (3.333,0.0175) (3.500,0.0132) (3.667,0.0098) (3.833,0.0072) (4.000,0.0052) (4.167,0.0037) (4.333,0.0026) (4.500,0.0018) (4.667,0.0012) (4.833,0.0008) (5.000,0.0005) (5.167,0.0004) (5.333,0.0002) (5.500,0.0001) (5.667,0.0001) (5.833,0.0001) (6.000,0.0000)};
\draw[cExact,very thick] plot coordinates {(-6.000,0.0006) (-5.833,0.0008) (-5.667,0.0011) (-5.500,0.0015) (-5.333,0.0020) (-5.167,0.0027) (-5.000,0.0036) (-4.833,0.0047) (-4.667,0.0061) (-4.500,0.0079) (-4.333,0.0101) (-4.167,0.0128) (-4.000,0.0160) (-3.833,0.0199) (-3.667,0.0245) (-3.500,0.0299) (-3.333,0.0361) (-3.167,0.0433) (-3.000,0.0514) (-2.833,0.0604) (-2.667,0.0704) (-2.500,0.0813) (-2.333,0.0930) (-2.167,0.1053) (-2.000,0.1183) (-1.833,0.1315) (-1.667,0.1450) (-1.500,0.1583) (-1.333,0.1713) (-1.167,0.1836) (-1.000,0.1950) (-0.833,0.2052) (-0.667,0.2139) (-0.500,0.2209) (-0.333,0.2261) (-0.167,0.2293) (0.000,0.2303) (0.167,0.2293) (0.333,0.2261) (0.500,0.2209) (0.667,0.2139) (0.833,0.2052) (1.000,0.1950) (1.167,0.1836) (1.333,0.1713) (1.500,0.1583) (1.667,0.1450) (1.833,0.1315) (2.000,0.1183) (2.167,0.1053) (2.333,0.0930) (2.500,0.0813) (2.667,0.0704) (2.833,0.0604) (3.000,0.0514) (3.167,0.0433) (3.333,0.0361) (3.500,0.0299) (3.667,0.0245) (3.833,0.0199) (4.000,0.0160) (4.167,0.0128) (4.333,0.0101) (4.500,0.0079) (4.667,0.0061) (4.833,0.0047) (5.000,0.0036) (5.167,0.0027) (5.333,0.0020) (5.500,0.0015) (5.667,0.0011) (5.833,0.0008) (6.000,0.0006)};
\draw[cBar,densely dashed,thick] (2.0,0.440) -- (3.2,0.440) node[right,black] {變異數 1：0.3989};
\draw[cWalk,densely dashed,thick] (2.0,0.385) -- (3.2,0.385) node[right,black] {變異數 2：0.2821};
\draw[cExact,very thick] (2.0,0.330) -- (3.2,0.330) node[right,black] {變異數 3：0.2303};
\fill (-5,0.0036) circle (1.1pt);
\fill (-4,0.0160) circle (1.1pt);
\fill (-3,0.0514) circle (1.1pt);
\fill (-2,0.1183) circle (1.1pt);
\fill (-1,0.1950) circle (1.1pt);
\fill (0,0.2303) circle (1.1pt);
\fill (1,0.1950) circle (1.1pt);
\fill (2,0.1183) circle (1.1pt);
\fill (3,0.0514) circle (1.1pt);
\fill (4,0.0160) circle (1.1pt);
\fill (5,0.0036) circle (1.1pt);
\end{tikzpicture}}

\newcommand{\SevenTabSplit}{%
\begin{tabular}{@{}lll@{}}
\toprule
回原點 & 一維 & 二維\\
\midrule
4 步，$M=2$ & $1+4+1=6$ & $16+4+16=36$\\
4 步，$M=3$ & $2\cdot3=6$ & $4\cdot9=36$\\
6 步，$2+2+2$ & 20 & 400\\
暴力列舉 & 6，20 & 36，400\\
\bottomrule
\end{tabular}}

\newcommand{\SevenFigTwoStep}{%
\begin{tikzpicture}[x=0.17cm,y=0.17cm,font=\scriptsize]
\draw[gray!25] (-2.2,-2.2) rectangle (2.2,2.2);
\fill (0.0,0) circle (0.45);
\draw[cWalk,thick,-{Stealth[length=2.5pt]}] (0.0,0) -- (1.0,0) -- (2.0,0);
\draw[gray!25] (2.8,-2.2) rectangle (7.2,2.2);
\fill (5.0,0) circle (0.45);
\draw[cWalk,thick,-{Stealth[length=2.5pt]}] (5.0,0) -- (6.0,0) -- (6.0,1);
\draw[gray!25] (7.8,-2.2) rectangle (12.2,2.2);
\fill (10.0,0) circle (0.45);
\draw[gray,thick] (10.0,0) -- (11.0,0) -- (11.00,0.45) -- (10.00,0.45);
\node[cWrong,font=\small,inner sep=0pt] at (11.60,0.45) {$\times$};
\draw[gray!25] (12.8,-2.2) rectangle (17.2,2.2);
\fill (15.0,0) circle (0.45);
\draw[cWalk,thick,-{Stealth[length=2.5pt]}] (15.0,0) -- (16.0,0) -- (16.0,-1);
\node at (7.5,3.4) {第一步向東：3＋1};
\draw[gray!25] (20.4,-2.2) rectangle (24.8,2.2);
\fill (22.6,0) circle (0.45);
\draw[cWalk,thick,-{Stealth[length=2.5pt]}] (22.6,0) -- (22.6,1) -- (23.6,1);
\draw[gray!25] (25.4,-2.2) rectangle (29.8,2.2);
\fill (27.6,0) circle (0.45);
\draw[cWalk,thick,-{Stealth[length=2.5pt]}] (27.6,0) -- (27.6,1) -- (27.6,2);
\draw[gray!25] (30.4,-2.2) rectangle (34.8,2.2);
\fill (32.6,0) circle (0.45);
\draw[cWalk,thick,-{Stealth[length=2.5pt]}] (32.6,0) -- (32.6,1) -- (31.6,1);
\draw[gray!25] (35.4,-2.2) rectangle (39.8,2.2);
\fill (37.6,0) circle (0.45);
\draw[gray,thick] (37.6,0) -- (37.6,1) -- (37.15,1.00) -- (37.15,0.00);
\node[cWrong,font=\small,inner sep=0pt] at (37.15,1.60) {$\times$};
\node at (30.1,3.4) {向北：3＋1};
\draw[gray!25] (43.0,-2.2) rectangle (47.4,2.2);
\fill (45.2,0) circle (0.45);
\draw[gray,thick] (45.2,0) -- (44.2,0) -- (44.20,-0.45) -- (45.20,-0.45);
\node[cWrong,font=\small,inner sep=0pt] at (43.60,-0.45) {$\times$};
\draw[gray!25] (48.0,-2.2) rectangle (52.4,2.2);
\fill (50.2,0) circle (0.45);
\draw[cWalk,thick,-{Stealth[length=2.5pt]}] (50.2,0) -- (49.2,0) -- (49.2,1);
\draw[gray!25] (53.0,-2.2) rectangle (57.4,2.2);
\fill (55.2,0) circle (0.45);
\draw[cWalk,thick,-{Stealth[length=2.5pt]}] (55.2,0) -- (54.2,0) -- (53.2,0);
\draw[gray!25] (58.0,-2.2) rectangle (62.4,2.2);
\fill (60.2,0) circle (0.45);
\draw[cWalk,thick,-{Stealth[length=2.5pt]}] (60.2,0) -- (59.2,0) -- (59.2,-1);
\node at (52.7,3.4) {向西：3＋1};
\draw[gray!25] (65.6,-2.2) rectangle (70.0,2.2);
\fill (67.8,0) circle (0.45);
\draw[cWalk,thick,-{Stealth[length=2.5pt]}] (67.8,0) -- (67.8,-1) -- (68.8,-1);
\draw[gray!25] (70.6,-2.2) rectangle (75.0,2.2);
\fill (72.8,0) circle (0.45);
\draw[gray,thick] (72.8,0) -- (72.8,-1) -- (73.25,-1.00) -- (73.25,0.00);
\node[cWrong,font=\small,inner sep=0pt] at (73.25,-1.60) {$\times$};
\draw[gray!25] (75.6,-2.2) rectangle (80.0,2.2);
\fill (77.8,0) circle (0.45);
\draw[cWalk,thick,-{Stealth[length=2.5pt]}] (77.8,0) -- (77.8,-1) -- (76.8,-1);
\draw[gray!25] (80.6,-2.2) rectangle (85.0,2.2);
\fill (82.8,0) circle (0.45);
\draw[cWalk,thick,-{Stealth[length=2.5pt]}] (82.8,0) -- (82.8,-1) -- (82.8,-2);
\node at (75.3,3.4) {向南：3＋1};
\end{tikzpicture}}

\newcommand{\SevenTabSaw}{%
\begin{tabular}{@{}l*{7}{r}@{}}
\toprule
$N$ & 1 & 2 & 3 & 4 & 5 & 6 & 7\\
\midrule
$c_N$ & 4 & 12 & 36 & 100 & 284 & 780 & 2172\\
$c_N/4^N$ & 1.0000 & 0.7500 & 0.5625 & 0.3906 & 0.2773 & 0.1904 & 0.1326\\
$\langle R^2\rangle$ & 1.0000 & 2.6667 & 4.5556 & 7.0400 & 9.5634 & 12.5744 & 15.5562\\
\midrule
$N$ & 8 & 9 & 10 & 11 & 12 & 13 & 14\\
\midrule
$c_N$ & 5916 & 16268 & 44100 & 120292 & 324932 & 881500 & 2374444\\
$c_N/4^N$ & 0.0903 & 0.0621 & 0.0421 & 0.0287 & 0.0194 & 0.0131 & 0.0088\\
$\langle R^2\rangle$ & 19.0128 & 22.4114 & 26.2425 & 30.0177 & 34.1870 & 38.3043 & 42.7864\\
\bottomrule
\end{tabular}}

\newcommand{\SevenFigLogLog}{%
\begin{tikzpicture}[x=1.0cm,y=0.6cm,>=Stealth,font=\scriptsize]
\draw[->] (0,0) -- (4.3,0) node[right] {$N$};
\draw[->] (0,0) -- (0,6.4) node[above] {$\langle R^2\rangle$};
\foreach \x/\l in {0/1,1/2,2/4,3/8,4/16} {\draw (\x,0) -- (\x,-0.12); \node[below] at (\x,-0.1) {$\l$};}
\foreach \y/\l in {0/1,2/4,4/16,6/64} {\draw (0,\y) -- (-0.06,\y); \node[left] at (-0.06,\y) {$\l$};}
\draw[gray,densely dashed] (0,0) -- (4,4) node[below=1pt] {普通漫步 $N$};
\draw[cWalk,densely dotted,thick] (1.0,1.208) -- (4.0,5.708) node[above] {斜率 $\frac32$};
\fill[cBar] (0.000,0.000) circle (1.4pt);
\fill[cBar] (1.000,1.415) circle (1.4pt);
\fill[cBar] (1.585,2.188) circle (1.4pt);
\fill[cBar] (2.000,2.816) circle (1.4pt);
\fill[cBar] (2.322,3.258) circle (1.4pt);
\fill[cBar] (2.585,3.652) circle (1.4pt);
\fill[cBar] (2.807,3.959) circle (1.4pt);
\fill[cBar] (3.000,4.249) circle (1.4pt);
\fill[cBar] (3.170,4.486) circle (1.4pt);
\fill[cBar] (3.322,4.714) circle (1.4pt);
\fill[cBar] (3.459,4.908) circle (1.4pt);
\fill[cBar] (3.585,5.095) circle (1.4pt);
\fill[cBar] (3.700,5.259) circle (1.4pt);
\fill[cBar] (3.807,5.419) circle (1.4pt);
\node[cBar,left] at (3.707,5.519) {自迴避 42.79};
\end{tikzpicture}}

\newcommand{\SevenFigChords}{%
\begin{tikzpicture}[font=\scriptsize]
\draw[thick] (0.000,0) -- (0.990,0);
\draw[cWrong,densely dashed,thick] (0.330,0) arc[start angle=180,end angle=0,radius=0.165];
\fill (0.330,0) circle (1.2pt);
\fill (0.660,0) circle (1.2pt);
\node[below] at (0.495,-0.05) {一階 (7.19)};
\draw[thick] (1.440,0) -- (3.090,0);
\draw[cWrong,densely dashed,thick] (1.770,0) arc[start angle=180,end angle=0,radius=0.495];
\draw[cWrong,densely dashed,thick] (2.100,0) arc[start angle=180,end angle=0,radius=0.165];
\fill (1.770,0) circle (1.2pt);
\fill (2.100,0) circle (1.2pt);
\fill (2.430,0) circle (1.2pt);
\fill (2.760,0) circle (1.2pt);
\node[below] at (2.265,-0.05) {(a) 巢狀};
\draw[thick] (3.540,0) -- (5.190,0);
\draw[cWrong,densely dashed,thick] (3.870,0) arc[start angle=180,end angle=0,radius=0.165];
\draw[cWrong,densely dashed,thick] (4.530,0) arc[start angle=180,end angle=0,radius=0.165];
\fill (3.870,0) circle (1.2pt);
\fill (4.200,0) circle (1.2pt);
\fill (4.530,0) circle (1.2pt);
\fill (4.860,0) circle (1.2pt);
\node[below] at (4.365,-0.05) {(b) 並排};
\draw[thick] (5.640,0) -- (7.290,0);
\draw[cWrong,densely dashed,thick] (5.970,0) arc[start angle=180,end angle=0,radius=0.330];
\draw[cWrong,densely dashed,thick] (6.300,0) arc[start angle=180,end angle=0,radius=0.330];
\fill (5.970,0) circle (1.2pt);
\fill (6.300,0) circle (1.2pt);
\fill (6.630,0) circle (1.2pt);
\fill (6.960,0) circle (1.2pt);
\node[below] at (6.465,-0.05) {(c) 交叉};
\end{tikzpicture}}

\newcommand{\SevenFigThird}{%
\begin{tikzpicture}[font=\tiny]
\begin{scope}[yshift=0.00cm]
\draw[thick] (0.000,0) -- (1.260,0);
\draw[gray,densely dashed,thick] (0.180,0) arc[start angle=180,end angle=0,radius=0.090];
\draw[gray,densely dashed,thick] (0.540,0) arc[start angle=180,end angle=0,radius=0.090];
\draw[gray,densely dashed,thick] (0.900,0) arc[start angle=180,end angle=0,radius=0.090];
\fill (0.180,0) circle (0.9pt);
\fill (0.360,0) circle (0.9pt);
\fill (0.540,0) circle (0.9pt);
\fill (0.720,0) circle (0.9pt);
\fill (0.900,0) circle (0.9pt);
\fill (1.080,0) circle (0.9pt);
\end{scope}
\begin{scope}[yshift=0.00cm]
\draw[thick] (1.600,0) -- (2.860,0);
\draw[gray,densely dashed,thick] (1.780,0) arc[start angle=180,end angle=0,radius=0.090];
\draw[gray,densely dashed,thick] (2.140,0) arc[start angle=180,end angle=0,radius=0.180];
\draw[gray,densely dashed,thick] (2.320,0) arc[start angle=180,end angle=0,radius=0.180];
\fill (1.780,0) circle (0.9pt);
\fill (1.960,0) circle (0.9pt);
\fill (2.140,0) circle (0.9pt);
\fill (2.320,0) circle (0.9pt);
\fill (2.500,0) circle (0.9pt);
\fill (2.680,0) circle (0.9pt);
\end{scope}
\begin{scope}[yshift=0.00cm]
\draw[thick] (3.200,0) -- (4.460,0);
\draw[gray,densely dashed,thick] (3.380,0) arc[start angle=180,end angle=0,radius=0.090];
\draw[gray,densely dashed,thick] (3.740,0) arc[start angle=180,end angle=0,radius=0.270];
\draw[gray,densely dashed,thick] (3.920,0) arc[start angle=180,end angle=0,radius=0.090];
\fill (3.380,0) circle (0.9pt);
\fill (3.560,0) circle (0.9pt);
\fill (3.740,0) circle (0.9pt);
\fill (3.920,0) circle (0.9pt);
\fill (4.100,0) circle (0.9pt);
\fill (4.280,0) circle (0.9pt);
\end{scope}
\begin{scope}[yshift=0.00cm]
\draw[thick] (4.800,0) -- (6.060,0);
\draw[gray,densely dashed,thick] (4.980,0) arc[start angle=180,end angle=0,radius=0.180];
\draw[gray,densely dashed,thick] (5.160,0) arc[start angle=180,end angle=0,radius=0.180];
\draw[gray,densely dashed,thick] (5.700,0) arc[start angle=180,end angle=0,radius=0.090];
\fill (4.980,0) circle (0.9pt);
\fill (5.160,0) circle (0.9pt);
\fill (5.340,0) circle (0.9pt);
\fill (5.520,0) circle (0.9pt);
\fill (5.700,0) circle (0.9pt);
\fill (5.880,0) circle (0.9pt);
\end{scope}
\begin{scope}[yshift=0.00cm]
\draw[thick] (6.400,0) -- (7.660,0);
\draw[cWalk,densely dashed,thick] (6.580,0) arc[start angle=180,end angle=0,radius=0.180];
\draw[cWalk,densely dashed,thick] (6.760,0) arc[start angle=180,end angle=0,radius=0.270];
\draw[cWalk,densely dashed,thick] (7.120,0) arc[start angle=180,end angle=0,radius=0.180];
\fill (6.580,0) circle (0.9pt);
\fill (6.760,0) circle (0.9pt);
\fill (6.940,0) circle (0.9pt);
\fill (7.120,0) circle (0.9pt);
\fill (7.300,0) circle (0.9pt);
\fill (7.480,0) circle (0.9pt);
\end{scope}
\begin{scope}[yshift=-0.85cm]
\draw[thick] (0.000,0) -- (1.260,0);
\draw[cWalk,densely dashed,thick] (0.180,0) arc[start angle=180,end angle=0,radius=0.180];
\draw[cWalk,densely dashed,thick] (0.360,0) arc[start angle=180,end angle=0,radius=0.360];
\draw[cWalk,densely dashed,thick] (0.720,0) arc[start angle=180,end angle=0,radius=0.090];
\fill (0.180,0) circle (0.9pt);
\fill (0.360,0) circle (0.9pt);
\fill (0.540,0) circle (0.9pt);
\fill (0.720,0) circle (0.9pt);
\fill (0.900,0) circle (0.9pt);
\fill (1.080,0) circle (0.9pt);
\end{scope}
\begin{scope}[yshift=-0.85cm]
\draw[thick] (1.600,0) -- (2.860,0);
\draw[gray,densely dashed,thick] (1.780,0) arc[start angle=180,end angle=0,radius=0.270];
\draw[gray,densely dashed,thick] (1.960,0) arc[start angle=180,end angle=0,radius=0.090];
\draw[gray,densely dashed,thick] (2.500,0) arc[start angle=180,end angle=0,radius=0.090];
\fill (1.780,0) circle (0.9pt);
\fill (1.960,0) circle (0.9pt);
\fill (2.140,0) circle (0.9pt);
\fill (2.320,0) circle (0.9pt);
\fill (2.500,0) circle (0.9pt);
\fill (2.680,0) circle (0.9pt);
\end{scope}
\begin{scope}[yshift=-0.85cm]
\draw[thick] (3.200,0) -- (4.460,0);
\draw[cWalk,densely dashed,thick] (3.380,0) arc[start angle=180,end angle=0,radius=0.270];
\draw[cWalk,densely dashed,thick] (3.560,0) arc[start angle=180,end angle=0,radius=0.270];
\draw[cWalk,densely dashed,thick] (3.740,0) arc[start angle=180,end angle=0,radius=0.270];
\fill (3.380,0) circle (0.9pt);
\fill (3.560,0) circle (0.9pt);
\fill (3.740,0) circle (0.9pt);
\fill (3.920,0) circle (0.9pt);
\fill (4.100,0) circle (0.9pt);
\fill (4.280,0) circle (0.9pt);
\end{scope}
\begin{scope}[yshift=-0.85cm]
\draw[thick] (4.800,0) -- (6.060,0);
\draw[cWalk,densely dashed,thick] (4.980,0) arc[start angle=180,end angle=0,radius=0.270];
\draw[cWalk,densely dashed,thick] (5.160,0) arc[start angle=180,end angle=0,radius=0.360];
\draw[cWalk,densely dashed,thick] (5.340,0) arc[start angle=180,end angle=0,radius=0.180];
\fill (4.980,0) circle (0.9pt);
\fill (5.160,0) circle (0.9pt);
\fill (5.340,0) circle (0.9pt);
\fill (5.520,0) circle (0.9pt);
\fill (5.700,0) circle (0.9pt);
\fill (5.880,0) circle (0.9pt);
\end{scope}
\begin{scope}[yshift=-0.85cm]
\draw[thick] (6.400,0) -- (7.660,0);
\draw[cWalk,densely dashed,thick] (6.580,0) arc[start angle=180,end angle=0,radius=0.360];
\draw[cWalk,densely dashed,thick] (6.760,0) arc[start angle=180,end angle=0,radius=0.090];
\draw[cWalk,densely dashed,thick] (7.120,0) arc[start angle=180,end angle=0,radius=0.180];
\fill (6.580,0) circle (0.9pt);
\fill (6.760,0) circle (0.9pt);
\fill (6.940,0) circle (0.9pt);
\fill (7.120,0) circle (0.9pt);
\fill (7.300,0) circle (0.9pt);
\fill (7.480,0) circle (0.9pt);
\end{scope}
\begin{scope}[yshift=-1.70cm]
\draw[thick] (0.000,0) -- (1.260,0);
\draw[cWalk,densely dashed,thick] (0.180,0) arc[start angle=180,end angle=0,radius=0.360];
\draw[cWalk,densely dashed,thick] (0.360,0) arc[start angle=180,end angle=0,radius=0.180];
\draw[cWalk,densely dashed,thick] (0.540,0) arc[start angle=180,end angle=0,radius=0.270];
\fill (0.180,0) circle (0.9pt);
\fill (0.360,0) circle (0.9pt);
\fill (0.540,0) circle (0.9pt);
\fill (0.720,0) circle (0.9pt);
\fill (0.900,0) circle (0.9pt);
\fill (1.080,0) circle (0.9pt);
\end{scope}
\begin{scope}[yshift=-1.70cm]
\draw[thick] (1.600,0) -- (2.860,0);
\draw[cWalk,densely dashed,thick] (1.780,0) arc[start angle=180,end angle=0,radius=0.360];
\draw[cWalk,densely dashed,thick] (1.960,0) arc[start angle=180,end angle=0,radius=0.360];
\draw[cWalk,densely dashed,thick] (2.140,0) arc[start angle=180,end angle=0,radius=0.090];
\fill (1.780,0) circle (0.9pt);
\fill (1.960,0) circle (0.9pt);
\fill (2.140,0) circle (0.9pt);
\fill (2.320,0) circle (0.9pt);
\fill (2.500,0) circle (0.9pt);
\fill (2.680,0) circle (0.9pt);
\end{scope}
\begin{scope}[yshift=-1.70cm]
\draw[thick] (3.200,0) -- (4.460,0);
\draw[cWalk,densely dashed,thick] (3.380,0) arc[start angle=180,end angle=0,radius=0.450];
\draw[cWalk,densely dashed,thick] (3.560,0) arc[start angle=180,end angle=0,radius=0.090];
\draw[cWalk,densely dashed,thick] (3.920,0) arc[start angle=180,end angle=0,radius=0.090];
\fill (3.380,0) circle (0.9pt);
\fill (3.560,0) circle (0.9pt);
\fill (3.740,0) circle (0.9pt);
\fill (3.920,0) circle (0.9pt);
\fill (4.100,0) circle (0.9pt);
\fill (4.280,0) circle (0.9pt);
\end{scope}
\begin{scope}[yshift=-1.70cm]
\draw[thick] (4.800,0) -- (6.060,0);
\draw[cWalk,densely dashed,thick] (4.980,0) arc[start angle=180,end angle=0,radius=0.450];
\draw[cWalk,densely dashed,thick] (5.160,0) arc[start angle=180,end angle=0,radius=0.180];
\draw[cWalk,densely dashed,thick] (5.340,0) arc[start angle=180,end angle=0,radius=0.180];
\fill (4.980,0) circle (0.9pt);
\fill (5.160,0) circle (0.9pt);
\fill (5.340,0) circle (0.9pt);
\fill (5.520,0) circle (0.9pt);
\fill (5.700,0) circle (0.9pt);
\fill (5.880,0) circle (0.9pt);
\end{scope}
\begin{scope}[yshift=-1.70cm]
\draw[thick] (6.400,0) -- (7.660,0);
\draw[cWalk,densely dashed,thick] (6.580,0) arc[start angle=180,end angle=0,radius=0.450];
\draw[cWalk,densely dashed,thick] (6.760,0) arc[start angle=180,end angle=0,radius=0.270];
\draw[cWalk,densely dashed,thick] (6.940,0) arc[start angle=180,end angle=0,radius=0.090];
\fill (6.580,0) circle (0.9pt);
\fill (6.760,0) circle (0.9pt);
\fill (6.940,0) circle (0.9pt);
\fill (7.120,0) circle (0.9pt);
\fill (7.300,0) circle (0.9pt);
\fill (7.480,0) circle (0.9pt);
\end{scope}
\end{tikzpicture}}

\newcommand{\SevenTabIE}{%
\begin{tabular}{@{}l*{9}{r}@{}}
\toprule
$K$ & 0 & 1 & 2 & 3 & 4 & 5 & 6 & 7 & 8\\
\midrule
$T_K$（$N=4$） & $256$ & $228$ & $96$ & $28$ & $4$ &  &  &  & \\
$P_K$（$N=4$） & $256$ & $28$ & $124$ & $96$ & $\mathbf{100}$ &  &  &  & \\
$P_K$（$N=6$） & $4096$ & $-3152$ & $4116$ & $-1492$ & $1996$ & $268$ & $928$ & $748$ & $784$\\
$P_K$（$N=8$） & $65536$ & $-119332$ & $190940$ & $-222036$ & $248468$ & $-219088$ & $186604$ & $-119052$ & $79824$\\
\midrule
$K$ & 9 & 10 & 11 & 12 & 13 & 14 & 15 & 16 & \\
\midrule
$P_K$（$N=6$） & $\mathbf{780}$ &  &  &  &  &  &  &  & \\
$P_K$（$N=8$） & $-31528$ & $22296$ & $-216$ & $7784$ & $5496$ & $5976$ & $5912$ & $\mathbf{5916}$ & \\
\bottomrule
\end{tabular}}

\newcommand{\SevenFigDyson}{%
\begin{tikzpicture}[font=\small,baseline=(current bounding box.center)]
\node[anchor=east] at (0,0) {$g=$};
\draw[thick] (0.10,0) -- (0.65,0);
\node at (0.95,0) {$+$};
\draw[thick] (1.25,0) -- (1.80,0);
\draw[cWrong,thick,fill=cWrong!8] (2.02,0) circle (0.22) node {\scriptsize$\Sigma$};
\draw[thick] (2.24,0) -- (2.79,0);
\node at (3.09,0) {$+$};
\draw[thick] (3.39,0) -- (3.94,0);
\draw[cWrong,thick,fill=cWrong!8] (4.16,0) circle (0.22) node {\scriptsize$\Sigma$};
\draw[thick] (4.38,0) -- (4.93,0);
\draw[cWrong,thick,fill=cWrong!8] (5.15,0) circle (0.22) node {\scriptsize$\Sigma$};
\draw[thick] (5.37,0) -- (5.92,0);
\node at (6.22,0) {$+$};
\node[anchor=west] at (6.32,0) {$\cdots=\dfrac{g_0}{1-\Sigma g_0}=\dfrac{1}{g_0^{-1}-\Sigma}$};
\end{tikzpicture}}

\newcommand{\SevenFigDysonBars}{%
\begin{tikzpicture}[font=\scriptsize,baseline=(current bounding box.center)]
\begin{scope}
\draw[cBar,fill=cBar!25] (0.00,0) rectangle (0.36,0.720);
\node[below] at (0.18,0) {$2$};
\draw[cBar,fill=cBar!25] (0.48,0) rectangle (0.84,1.080);
\node[below] at (0.66,0) {$3$};
\draw[cBar,fill=cBar!25] (0.96,0) rectangle (1.32,1.260);
\node[below] at (1.14,0) {$\frac{7}{2}$};
\draw[cBar,fill=cBar!25] (1.44,0) rectangle (1.80,1.350);
\node[below] at (1.62,0) {$\frac{15}{4}$};
\draw[cBar,fill=cBar!25] (1.92,0) rectangle (2.28,1.395);
\node[below] at (2.10,0) {$\frac{31}{8}$};
\draw[cExact,densely dashed,thick] (-0.1,1.440) -- (2.40,1.440) node[right] {$4$};
\draw (-0.1,0) -- (2.40,0);
\node[below] at (1.20,-0.5) {$g_0=2,\ \Sigma=\frac14$ 的部分和};
\end{scope}
\end{tikzpicture}}

\newcommand{\SevenTabDyson}{%
\begin{tabular}[c]{@{}lrrrrrl@{}}
\toprule
摺積 (7.20)：$N$ & 0 & 1 & 2 & 3 & 4 & \\
\quad $H_1=H_2=2^K$ & 1 & 2 & 4 & 8 & 16 & \\
\quad $\sum_MH_1(M)H_2(N-M)$ & 1 & 4 & 12 & 32 & 80 & $z=0.1$：1.5625\\
\midrule
Dyson：項數 $n$ & 1 & 2 & 3 & 4 & 5 & $\infty$\\
\quad $g_0\sum_{k<n}(\Sigma g_0)^k$ & $2$ & $3$ & $\frac{7}{2}$ & $\frac{15}{4}$ & $\frac{31}{8}$ & $4$\\
\bottomrule
\end{tabular}
\qquad
\begin{tabular}[c]{@{}crrr@{}}
\toprule
計數版 $k$ & $(2k-1)!!$ & 不可約 $I_k$ & $\frac1{1-I(x)}$\\
\midrule
0 & 1 & --- & 1\\
1 & 1 & 1 & 1\\
2 & 3 & 2 & 3\\
3 & 15 & 10 & 15\\
4 & 105 & 74 & 105\\
\bottomrule
\end{tabular}}

\newcommand{\SevenTabPower}{%
\begin{tabular}{@{}c|rrrrr@{}}
\toprule
$d$ \textbackslash\ $l$ & 0 & 1 & 2 & 3 & 4\\
\midrule
1 & $1$ & $-\frac{1}{2}$ & $-2$ & $-\frac{7}{2}$ & $-5$\\
2 & $1$ & $0$ & $-1$ & $-2$ & $-3$\\
3 & $1$ & $\frac{1}{2}$ & $0$ & $-\frac{1}{2}$ & $-1$\\
4 & $1$ & $1$ & $1$ & $1$ & $1$\\
5 & $1$ & $\frac{3}{2}$ & $2$ & $\frac{5}{2}$ & $3$\\
\bottomrule
\end{tabular}}

\newcommand{\SevenTabMF}{%
\begin{tabular}{@{}lccccc@{}}
\toprule
練習 7.4：$(m,\mathcal N)$ & $(10,10^2)$ & $(10^2,10^4)$ & $(10^3,10^6)$ & & 極限\\
\quad $\mathcal Nz_c^mC_{MF}(m)$ & 0.3385 & 0.3643 & 0.3675 & & 0.3679\\
\midrule
練習 7.5：$\mathcal N$ & $10$ & $10^2$ & $10^3$ & $10^4$ & \\
\quad $\frac z{z_c}=\frac12$：$\mathcal N(1-\frac z{z_c})G_{MF}$ & 0.6899 & 0.9341 & 0.9922 & 0.9992 & 1\\
\quad $\frac z{z_c}=\frac32$：$\frac{\ln G_{MF}}{\mathcal N}$ & --- & 0.0452 & 0.0596 & 0.0621 & 0.0625\\
\bottomrule
\end{tabular}}

\newcommand{\SevenFigFlory}{%
\begin{tikzpicture}[x=0.0085cm,y=0.022cm,>=Stealth,font=\scriptsize]
\draw[->] (0,0) -- (640,0) node[right] {$R$};
\draw[->] (0,-160) -- (0,12);
\foreach \x in {100,200,300,400,500,600} {\draw (\x,0) -- (\x,3); \node[above] at (\x,2) {\x};}
\foreach \y in {-150,-100,-50} {\draw (0,\y) -- (-8,\y); \node[left] at (-8,\y) {$\y$};}
\begin{scope}\clip (0,-160) rectangle (640,0);
\draw[cBar,thick] plot coordinates {(50,-2.50) (55,-3.02) (60,-3.60) (65,-4.22) (70,-4.90) (75,-5.62) (80,-6.40) (85,-7.22) (90,-8.10) (95,-9.03) (100,-10.00) (105,-11.03) (110,-12.10) (115,-13.22) (120,-14.40) (125,-15.62) (130,-16.90) (135,-18.23) (140,-19.60) (145,-21.02) (150,-22.50) (155,-24.02) (160,-25.60) (165,-27.23) (170,-28.90) (175,-30.62) (180,-32.40) (185,-34.23) (190,-36.10) (195,-38.02) (200,-40.00) (205,-42.02) (210,-44.10) (215,-46.23) (220,-48.40) (225,-50.62) (230,-52.90) (235,-55.23) (240,-57.60) (245,-60.02) (250,-62.50) (255,-65.03) (260,-67.60) (265,-70.22) (270,-72.90) (275,-75.62) (280,-78.40) (285,-81.22) (290,-84.10) (295,-87.03) (300,-90.00) (305,-93.03) (310,-96.10) (315,-99.22) (320,-102.40) (325,-105.62) (330,-108.90) (335,-112.22) (340,-115.60) (345,-119.03) (350,-122.50) (355,-126.03) (360,-129.60) (365,-133.22) (370,-136.90) (375,-140.62) (380,-144.40) (385,-148.22) (390,-152.10) (395,-156.03) (400,-160.00) (405,-164.03) (410,-168.10) (415,-172.22) (420,-176.40) (425,-180.62) (430,-184.90) (435,-189.22) (440,-193.60) (445,-198.03) (450,-202.50) (455,-207.03) (460,-211.60) (465,-216.22) (470,-220.90) (475,-225.62) (480,-230.40) (485,-235.22) (490,-240.10) (495,-245.03) (500,-250.00) (505,-255.03) (510,-260.10) (515,-265.23) (520,-270.40) (525,-275.62) (530,-280.90) (535,-286.23) (540,-291.60) (545,-297.02) (550,-302.50) (555,-308.02) (560,-313.60) (565,-319.23) (570,-324.90) (575,-330.62) (580,-336.40) (585,-342.23) (590,-348.10) (595,-354.02) (600,-360.00)};
\draw[cWalk,thick] plot coordinates {(50,-400.00) (55,-330.58) (60,-277.78) (65,-236.69) (70,-204.08) (75,-177.78) (80,-156.25) (85,-138.41) (90,-123.46) (95,-110.80) (100,-100.00) (105,-90.70) (110,-82.64) (115,-75.61) (120,-69.44) (125,-64.00) (130,-59.17) (135,-54.87) (140,-51.02) (145,-47.56) (150,-44.44) (155,-41.62) (160,-39.06) (165,-36.73) (170,-34.60) (175,-32.65) (180,-30.86) (185,-29.22) (190,-27.70) (195,-26.30) (200,-25.00) (205,-23.80) (210,-22.68) (215,-21.63) (220,-20.66) (225,-19.75) (230,-18.90) (235,-18.11) (240,-17.36) (245,-16.66) (250,-16.00) (255,-15.38) (260,-14.79) (265,-14.24) (270,-13.72) (275,-13.22) (280,-12.76) (285,-12.31) (290,-11.89) (295,-11.49) (300,-11.11) (305,-10.75) (310,-10.41) (315,-10.08) (320,-9.77) (325,-9.47) (330,-9.18) (335,-8.91) (340,-8.65) (345,-8.40) (350,-8.16) (355,-7.93) (360,-7.72) (365,-7.51) (370,-7.30) (375,-7.11) (380,-6.93) (385,-6.75) (390,-6.57) (395,-6.41) (400,-6.25) (405,-6.10) (410,-5.95) (415,-5.81) (420,-5.67) (425,-5.54) (430,-5.41) (435,-5.28) (440,-5.17) (445,-5.05) (450,-4.94) (455,-4.83) (460,-4.73) (465,-4.62) (470,-4.53) (475,-4.43) (480,-4.34) (485,-4.25) (490,-4.16) (495,-4.08) (500,-4.00) (505,-3.92) (510,-3.84) (515,-3.77) (520,-3.70) (525,-3.63) (530,-3.56) (535,-3.49) (540,-3.43) (545,-3.37) (550,-3.31) (555,-3.25) (560,-3.19) (565,-3.13) (570,-3.08) (575,-3.02) (580,-2.97) (585,-2.92) (590,-2.87) (595,-2.82) (600,-2.78)};
\draw[cExact,very thick] plot coordinates {(50,-400.00) (55,-333.60) (60,-281.38) (65,-240.91) (70,-208.98) (75,-183.40) (80,-162.65) (85,-145.63) (90,-131.56) (95,-119.83) (100,-110.00) (105,-101.73) (110,-94.74) (115,-88.84) (120,-83.84) (125,-79.62) (130,-76.07) (135,-73.09) (140,-70.62) (145,-68.59) (150,-66.94) (155,-65.65) (160,-64.66) (165,-63.96) (170,-63.50) (175,-63.28) (180,-63.26) (185,-63.44) (190,-63.80) (195,-64.32) (200,-65.00) (205,-65.82) (210,-66.78) (215,-67.86) (220,-69.06) (225,-70.38) (230,-71.80) (235,-73.33) (240,-74.96) (245,-76.68) (250,-78.50) (255,-80.40) (260,-82.39) (265,-84.46) (270,-86.62) (275,-88.85) (280,-91.16) (285,-93.54) (290,-95.99) (295,-98.52) (300,-101.11) (305,-103.77) (310,-106.51) (315,-109.30) (320,-112.17) (325,-115.09) (330,-118.08) (335,-121.14) (340,-124.25) (345,-127.43) (350,-130.66) (355,-133.96) (360,-137.32) (365,-140.73) (370,-144.20) (375,-147.74) (380,-151.33) (385,-154.97) (390,-158.67) (395,-162.43) (400,-166.25) (405,-170.12) (410,-174.05) (415,-178.03) (420,-182.07) (425,-186.16) (430,-190.31) (435,-194.51) (440,-198.77) (445,-203.07) (450,-207.44) (455,-211.86) (460,-216.33) (465,-220.85) (470,-225.43) (475,-230.06) (480,-234.74) (485,-239.48) (490,-244.26) (495,-249.11) (500,-254.00) (505,-258.95) (510,-263.94) (515,-269.00) (520,-274.10) (525,-279.25) (530,-284.46) (535,-289.72) (540,-295.03) (545,-300.39) (550,-305.81) (555,-311.27) (560,-316.79) (565,-322.36) (570,-327.98) (575,-333.65) (580,-339.37) (585,-345.15) (590,-350.97) (595,-356.85) (600,-362.78)};
\end{scope}
\draw[gray,densely dashed] (177.83,-160) -- (177.83,0);
\fill[cExact] (177.83,-63.25) circle (1.6pt);
\fill (177.83,-31.62) circle (1.4pt);
\node[below right] at (177.83,-150) {$R^*=177.83$};
\node[cWalk,anchor=west] at (320,-32) {排斥 $-um^2/R^2$};
\node[cBar,anchor=west] at (320,-49) {熵 $-R^2/m$};
\node[cExact,anchor=west] at (320,-66) {和，最大 $-63.25$};
\node[black,anchor=west] at (320,-83) {交點：兩項都是 $-31.62$};
\end{tikzpicture}}

\newcommand{\SevenTabFlory}{%
\begin{tabular}{@{}ccccl@{}}
\toprule
$d$ & $\nu_F=\frac3{d+2}$ & 數值斜率 & $R^*$（$m=10^3$） & 已知值\\
\midrule
1 & $1$ & 1.0000 & 793.70 & 1\\
2 & $\frac{3}{4}$ & 0.7500 & 177.83 & $\frac34$（精確）\\
3 & $\frac{3}{5}$ & 0.6000 & 68.43 & 0.588（數值）\\
4 & $\frac{1}{2}$ & 0.5000 & 35.50 & $\frac12$\\
\bottomrule
\end{tabular}
\\[6pt]
\begin{tabular}{@{}l*{5}{c}c@{}}
\toprule
二維 $N$ & 6 & 8 & 10 & 12 & 14 & $\infty$\\
\midrule
$2\nu_{\rm eff}$ & 1.4306 & 1.4372 & 1.4442 & 1.4505 & 1.4556 & $\frac32$\\
\bottomrule
\end{tabular}}

\newcommand{\SevenTabErrata}{%
\begin{tabular}{@{}lll@{}}
\toprule
位置 & 書上 & 應為\\
\midrule
p.~169--173，(7.6)(7.11)、Ex.~7.1 & {\color{cWrong}$(2\pi Ma^2)^{-3/2}$} & {\color{cExact}$(2\pi Ma^2/3)^{-3/2}$}\\
p.~173，(7.11) & {\color{cWrong}$e^{-3|\vec x-\vec r_1|^2/3M_1a^2}$、$\cdots/3M_ka^2$} & {\color{cExact}分母 $2M_1a^2$、$2M_ka^2$}\\
p.~176，(7.21) 第二、三行 & {\color{cWrong}$Z^{N-M_1}$、$Z^{K_2}$} & {\color{cExact}$z^{N-M_1}$、$z^{K_2}$}\\
p.~178，(7.27) 前 & {\color{cWrong}Figure 7.3(a)} & {\color{cExact}Figure 7.3(b)（並排）}\\
p.~183，(7.33) & {\color{cWrong}$g_0[\Sigma_1+\Sigma_2]^2g_0$} & {\color{cExact}$g_0[\Sigma_1+\Sigma_2]g_0[\Sigma_1+\Sigma_2]g_0$}\\
p.~183，(7.35) & {\color{cWrong}$\frac{g_0}{1+(\Sigma_1+\Sigma_2)g_0}$} & {\color{cExact}$\frac{g_0}{1-(\Sigma_1+\Sigma_2)g_0}$}\\
p.~188，(7.47) & {\color{cWrong}$za^2\nabla^2G-(z/z_c-1)G=a^d\delta$} & {\color{cExact}$(1-z/z_c)G-za^2\nabla^2G=a^d\delta$}\\
p.~188，(7.48) & {\color{cWrong}$\psi_i(\vec r)=e^{-\vec q\cdot\vec r}$} & {\color{cExact}$\psi_i(\vec r)=e^{i\vec q\cdot\vec r}/\sqrt V$}\\
p.~188，(7.50) 後 & {\color{cWrong}(7.43dd)} & {\color{cExact}(7.43d)}\\
p.~189，(7.55) 前 & {\color{cWrong}積分 $=1/(z_0-z)$} & {\color{cExact}積分 $=z_c/(z_0-z)$}\\
p.~190，(7.58) 與前一句 & {\color{cWrong}$e^{-\mathcal Nz\lambda/z_c}$；$z^m\cdots\lambda^n$} & {\color{cExact}$e^{+\mathcal Nz\lambda/z_c}$；$(z/z_c)^m\cdots\lambda^m$}\\
p.~190，(7.59) & {\color{cWrong}$e^{-\mathcal N(f\lambda^2+\mathcal N\lambda)+m\ln\lambda}$} & {\color{cExact}$e^{-\mathcal N(f\lambda^2+\lambda)+m\ln\lambda}$}\\
p.~191，(7.61) 左邊 & {\color{cWrong}$C_{MF}(N;\vec x,\vec y)$} & {\color{cExact}$C_{MF}(m;\vec x,\vec y)$}\\
\bottomrule
\end{tabular}}


\section{路徑積分與自迴避（第 7 章）}

\begin{keybox}[先看這裡]
前六章的醉漢沒有記憶：每一步都和以前無關。這一章先換個角度看漫步——\textbf{把所有路徑加起來}（路徑和、路徑積分）；
再讓醉漢\textbf{不准踩到自己走過的點}（自迴避漫步，高分子的最簡單模型）。有了記憶，問題就難了：書上先把自迴避\textbf{一階一階展開}
（$1-f$ 展開、費曼圖、Dyson 方程），發現四維以下低階的展開一定失效；再用\textbf{平均場}（只留一個模態）與 \textbf{Flory 的論證}
（熵與排斥拔河）得到 $R\propto m^{3/(d+2)}$：二維的 $\frac34$ 剛好是精確值，三維的 0.6 也很接近數值的 0.588。總整理在 7.4 節。

\textbf{記號}：$C(N;\vec x,\vec y)$＝從 $\vec x$ 走 $N$ 步到 $\vec y$ 的走法數，格距 $a$；一條 $N$ 步的漫步走過 $N+1$ 個點 $\vec w_0=\vec x,\dots,\vec w_N$，
「嚴格」自迴避＝這 $N+1$ 個點兩兩不同。$c_N$＝二維 $N$ 步自迴避漫步的條數，$\langle R^2\rangle$＝終點平方距離的平均（普通漫步 $=N$）。
逸度 $z$（$z_c=\frac1{2d}$）、傳播子 $g_0=\frac1{1-z\chi(\vec q)}$、自能 $\Sigma$。平均場與 Flory 的步數寫 $m$，格點數寫 $\mathcal N$（書上 §7.2.6 起也寫成 $N$）。

\textbf{全章的主例}：二維 4 步回原點 36 條；4 步自迴避 100 條＝部分和 $256\to28\to124\to96\to100$；Flory 的 $d=2$、$m=1000$：$R^*=177.83$。
\begin{center}
\small
\begin{tabular}{@{}lll@{\qquad}lll@{}}
\toprule
節 & 內容 & 動畫 & 節 & 內容 & 動畫\\
\midrule
7.1 & 路徑和、高斯、路徑積分 & W701 & 7.2.3--7.2.4 & 生成函數、費曼規則、Dyson & W704\\
7.2--7.2.1 & 自迴避漫步、精確計數 & W702 & 7.2.5--7.2.6 & 微擾失效、平均場 & W705\\
7.2.2 & $1-f$ 展開、弧線圖 & W703 & 7.2.7 & Flory 的論證、回顧 & W706\\
\bottomrule
\end{tabular}
\end{center}
（上表是書上的節號，本講義 7.1、7.2.1--7.2.7 與它一一對應。）數字一律由 \texttt{codes/utils/walk\_saw.py} 重算，圖表由 \texttt{ch7/make\_figs.py} 產生並先做數字檢查；
逐式校對見 \texttt{ref/md/ch07\_formulas.md}。連續播放：\texttt{Ch7Demo}（W701--W706）。
\end{keybox}

% =====================================================================
\subsection{漫步＝路徑的和}
\leavevmode\bk{pp.~167--171}\ \anim{W701}\par\nopagebreak

\textbf{問}：從原點走 4 步回到原點，一維有幾條？二維呢？先猜。

\textbf{路徑的和}（(7.1)）：走法數就是「每條路徑的權重加起來」，格子上每條路徑的權重都是 1。把 $N$ 步切成兩段：先走 $M$ 步到中間點 $\vec r$，
再走 $N-M$ 步到 $\vec y$，然後對中間點求和——像轉機：從台北到紐約的走法，就是「經過每一個轉機城市」的走法加起來：
\[C(N;\vec x,\vec y)=\sum_{\vec r}C(M;\vec x,\vec r)\,C(N-M;\vec r,\vec y).\tag{7.3}\]

\textbf{算到底}（圖 \ref{seven-fig-split}）：一維 $N=4$、$M=2$。走 2 步能到的中間點是 $r=-2,0,2$，進來的走法 $C(2;r)=1,2,1$，出去（從 $r$ 再走 2 步回 0）也是 $1,2,1$，
所以 $C(4;0,0)=1\cdot1+2\cdot2+1\cdot1=6=\binom42$。二維：2 步能到 9 個中間點。中心 $C(2;\vec 0)=4$（出去又回來，4 個方向）、四個 $(\pm2,0),(0,\pm2)$ 各 1 種、四個 $(\pm1,\pm1)$ 各 2 種（先橫後直或先直後橫）；
回程的走法和進來一樣，所以
\[C(4;\vec0,\vec0)=4^2+4\cdot1^2+4\cdot2^2=16+4+16=36=\tbinom42^2 .\]
\textbf{揭曉}：6 與 36；暴力列舉 $2^4$、$4^4$ 條也是這兩個數。切在哪裡都一樣：$M=N-1$ 就是第 2 章的遞迴 (2.7)
（一維 2 個鄰居、各 $C(3;\pm1)=3$ 種：$3+3=6$；二維 $4\cdot9=36$）。切成 $k$ 段 (7.4)(7.5) 就對 $k-1$ 個中間點求和：6 步回原點切 $2+2+2$，一維 $20=\binom63$、二維 $400=\binom63^2$。

\textbf{每段都很長：高斯}（(7.6)）。每段走很多步時換成第 2 章的高斯；三維簡單立方格子走 $M$ 步，一個直角分量的變異數是 $Ma^2/3$：
\[C(N;\vec x,\vec y)=\int\!\!\cdots\!\!\int d\vec r_1\cdots d\vec r_{k-1}\prod_{l=1}^{k}\Big(\frac{2\pi M_la^2}{3}\Big)^{-3/2}e^{-3|\vec r_{l-1}-\vec r_l|^2/2M_la^2},
\qquad\vec r_0=\vec x,\ \vec r_k=\vec y.\tag{7.6}\]
兩段高斯接起來還是高斯，而且\textbf{變異數相加}（練習 7.1）。例：$M_1=3$、$M_2=6$（$a=1$）：一個分量的變異數 1、2，峰高 0.3989、0.2821；
摺積的變異數 $1+2=3$、峰高 $0.2303=1/\sqrt{2\pi\cdot3}$（圖 \ref{seven-fig-split} 右，點是數值摺積）。
\begin{errbox}[p.~169--173，(7.6)(7.11)、練習 7.1]
前因子都寫成 $(2\pi Ma^2)^{-3/2}$，少了 3：應為 $(2\pi Ma^2/3)^{-3/2}$。照書上，每做一次摺積就多出 $3^{-3/2}=0.1925$（一個分量 $3^{-1/2}=0.5774$），練習 7.1 就證不出來。
〔注〕(7.6) 也省略了每段的 $6^{M_l}$（乘起來是 $6^N$）與「格點和 $\to$ 積分」的 $a^{-3}$；只看形狀時無妨。
\end{errbox}

\textbf{切成無限多段：路徑積分}（(7.7)(7.8)）。令 $\Delta s=\Delta N_i$，每段的指數 $-\frac{3}{2a^2}\frac{|\Delta\vec r_i|^2}{\Delta N_i}=-\frac{3}{2a^2}\big|\frac{\Delta\vec r_i}{\Delta s}\big|^2\Delta s$
像「速度平方」，全部加起來變成積分：
\[C(N;\vec x,\vec y)\propto\prod_s\int d\vec r(s)\,\exp\Big[-\frac{3}{2a^2}\int_0^N\Big|\frac{d\vec r}{ds}\Big|^2ds\Big],\qquad\vec r(0)=\vec x,\ \vec r(N)=\vec y.\tag{7.8}\]
這就是 Feynman 的路徑積分：$s$ 是「時間」，$|d\vec r/ds|^2$ 是「動能」。〔注〕$\Delta N_i$ 必須 $\gg1$ 高斯才成立，所以 $d\vec r/ds$ 只是記號；(7.8) 的比例常數是無限大（紫外發散）。

\begin{figure}[htbp]
\centering
\SevenFigGrid\hfill{\footnotesize\setlength{\tabcolsep}{4pt}\SevenTabSplit}\hfill\SevenFigGauss
\caption{路徑的和。左：二維 4 步回原點（起點＝終點＝中心），9 個中間點上寫 $C(2;\vec r)$，「進來 × 出去」$=C(2;\vec r)^2$。
中：切法不同、答案相同。右：練習 7.1，變異數 1、2 的高斯（虛線，數字是峰高）摺積成變異數 3（實線），黑點是數值摺積}\label{seven-fig-split}
\end{figure}

% =====================================================================
\subsection{自迴避漫步}
\leavevmode\bk{pp.~171--192}\quad 高分子是一串單體連成的長鏈，兩個單體不能佔同一個位置。最簡單的模型：格子上的漫步，不准踩到自己走過的點。
這一節的難處全在「漫步有了記憶」。

\subsubsection{自迴避漫步與交互作用}
\leavevmode\bk{pp.~171--172}\ \anim{W702}\par\nopagebreak

\textbf{問}：如果醉漢不准踩到自己走過的點，10 步之後離起點會比較遠還是比較近？

\textbf{三種自迴避}：嚴格（踩到走過的點就停，那條路徑不算）、溫和（越靠近越可能停）、「真」自迴避（偏向不踩，但一直走）。本章用嚴格的。
寫成路徑和，就是每條路徑乘上權重 $\prod_{i>j}e^{-v(|\vec w_i-\vec w_j|)}$（(7.9)）：
兩點重合時 $v=+\infty$、權重 0，那條路徑被刪掉；不重合時 $v=0$、權重 1。(7.10) 再乘上與起點、終點的因子 $e^{-v(|\vec w_i-\vec x|)-v(|\vec w_i-\vec y|)}$。

\textbf{數一數}（二維，圖 \ref{seven-fig-saw}）：2 步的 16 條中，4 條第二步原地折返，剩 12 條；3 步的 64 條剩 36 條（第一步向東的 9 條 $\times4$）；4 步的 256 條剩 100 條。
存活的比例 $c_N/4^N$ 一路掉，$N=14$ 只剩 0.0088（$c_{14}=2374444$）；每多一步條數大約乘 $c_{14}/c_{13}=2.6936$（連接常數 $\mu\approx2.638$，只引用），比普通漫步的 4 小。
\textbf{揭曉}：條數雖少，剩下的卻走得遠：$\langle R^2\rangle$ 在 $N=14$ 是 42.79，普通漫步只有 14。種子 702 抽一條 10 步的自迴避漫步，終點 $(-3,3)$、$R^2=18$；
同一個種子的普通漫步終點 $(1,-1)$、$R^2=2$，中間踩回走過的點 3 次。自迴避漫步被自己的足跡推著走：在對數圖上 $\langle R^2\rangle\propto N^{2\nu}$ 的斜率比 1 大，往 $\frac32$ 爬（7.2.7 節）。

\begin{figure}[htbp]
\centering
\SevenFigTwoStep\\[6pt]
\begin{minipage}[c]{0.6\linewidth}\centering\footnotesize\setlength{\tabcolsep}{2.5pt}\SevenTabSaw\end{minipage}\hfill
\begin{minipage}[c]{0.38\linewidth}\centering\SevenFigLogLog\end{minipage}
\caption{二維自迴避漫步。上：16 條 2 步漫步，按第一步分四組；橙色 12 條自迴避，灰色 4 條原地折返（$\times$）。
左下：精確列舉 $N=1..14$（條數、存活比例、平均平方距離）。右下：$\langle R^2\rangle$ 對 $N$ 的對數圖：藍點是自迴避，虛線是普通漫步，點線是斜率 $\frac32$ 的參考線}\label{seven-fig-saw}
\end{figure}

\subsubsection{路徑積分與 \texorpdfstring{$1-f$}{1-f} 展開}
\leavevmode\bk{pp.~172--176}\ \anim{W703}\par\nopagebreak

\textbf{問}：「不准踩到自己」很難直接算。能不能先算普通漫步，再一點一點扣掉？

\textbf{連續版}（(7.11)--(7.16)）：把 (7.9) 的權重放進 (7.6)；交互作用取極短程的 $\delta$ 型，(7.8) 的指數多一項：
\[C(N;\vec x,\vec y)\propto\prod_s\int d\vec r(s)\exp\Big[-\frac{3}{2a^2}\int_0^N\Big|\frac{d\vec r}{ds}\Big|^2ds-v_0\int_0^N\!\!\int_0^{s_1}ds_1ds_2\,\delta\big(\vec r(s_1)-\vec r(s_2)\big)\Big].\tag{7.16}\]
〔注〕書上 (7.13)(7.14) 寫 $-v_0\int\!\!\int v(\cdots)$；$v_0$ 應只出現在 $\delta$ 的版本 (7.15)(7.16)。
\begin{errbox}[p.~173，(7.11)]
兩個指數印成 $-3|\vec x-\vec r_1|^2/3M_1a^2$、$-3|\vec r_k-\vec y|^2/3M_ka^2$；分母應為 $2M_1a^2$、$2M_ka^2$（同 (7.6)）。
照書上，一個分量的變異數變成 $Ma^2/2$（$M=6$：3），格子漫步卻是 $Ma^2/3$（2）。
\end{errbox}

\textbf{$1-f$ 展開}：令
\[e^{-v(x)}=1+\big(e^{-v(x)}-1\big)\equiv1-f(x),\qquad0\le f\le1,\tag{7.17}\]
$f$ 只在 $x\approx0$ 附近不是 0；嚴格自迴避時「重合」$f=1$、否則 $f=0$。每一對點乘一個 $(1-f)$（(7.18)），再乘開：0 階是普通漫步；
一階是「挑一對點、要它們重合」，扣掉：
\[-\sum_{M_1+M_2+M_3=N}\ \sum_{\vec u_1,\vec u_2}C(M_1;\vec x,\vec u_1)C(M_2;\vec u_1,\vec u_2)C(M_3;\vec u_2,\vec y)\,f(|\vec u_1-\vec u_2|),\tag{7.19}\]
畫成一條線加一條虛線弧（Fig.~7.2；圖 \ref{seven-fig-chord}）：走到 $\vec u_1$、繞一圈回到 $\vec u_2\approx\vec u_1$、再走到終點。
二階有 3 種（Fig.~7.3）：(a) 巢狀、(b) 並排、(c) 交叉；三階有 15 種（練習 7.2）。一般是 $2k$ 個點配成 $k$ 條弧，$(2k-1)!!=1,3,15,105$ 種——和第 6 章的 Wick 配對一樣。

\textbf{算到底}（二維 $N=4$）：嚴格自迴避的 $f$ 非 0 即 1，所以第 $k$ 階 $T_k$＝「每條漫步從它的重合點對中挑 $k$ 對」的方法數總和，而
$c_N=\sum_k(-1)^kT_k$——這就是排容原理。5 個點 $\vec w_0..\vec w_4$ 中可能重合的只有隔偶數步的：隔 2 步的 3 對 $(0,2),(1,3),(2,4)$，隔 4 步的 1 對 $(0,4)$。
一階：隔 2 步的一對要重合，中間 2 步得走回原處（$C(2;\vec0)=4$ 種），另外 2 步隨便走（$4^2$ 種）；隔 4 步的一對要 $C(4;\vec0)=36$ 種，所以
$T_1=3\cdot4^2\cdot4+36=192+36=228$。全部是 $T_k=256,228,96,28,4$，部分和
\[256\ \to\ 256-228=28\ \to\ 28+96=124\ \to\ 124-28=96\ \to\ 96+4=100=c_4 .\]
二階的 96 依兩對點的位置分：巢狀 16、交叉 32、共用一個點 48（格子上兩條弧可以共用端點；$N=4$ 沒有並排的）。
\textbf{揭曉}（圖 \ref{seven-fig-chord} 下）：$N=8$ 的部分和是 $65536\to-119332\to190940\to-222036\to\cdots\to5916$，共 17 項，前幾項連正負號都不對；
要全部加完才剛好是 $c_8=5916$。只取前幾階，數字可能錯得離譜——7.2.5 節會說明這不是偶然。

\begin{figure}[htbp]
\centering
\SevenFigChords\hfill\SevenFigThird\\[6pt]
{\footnotesize\setlength{\tabcolsep}{4pt}\SevenTabIE}
\caption{$1-f$ 展開。左上：實線是漫步、點是漫步上的位置、虛線弧＝要求兩點重合的 $-f$；一階一種、二階三種（書上 Fig.~7.2、7.3）。
右上：三階 15 種（練習 7.2），橙色 10 種不可約、灰色 5 種可約。下：二維嚴格自迴避的第 $K$ 階 $T_K$ 與部分和 $P_K=\sum_{k\le K}(-1)^kT_k$；粗體的最後一項就是 $c_N$}\label{seven-fig-chord}
\end{figure}

\subsubsection{生成函數的微擾展開與費曼規則}
\leavevmode\bk{pp.~176--182}\ \anim{W704}\par\nopagebreak

\textbf{問}：展開的每一項都是「切段再求和」的摺積，越高階越難算。有沒有辦法把摺積變成乘法？

\textbf{摺積變乘積}（(7.20)--(7.22)）：令 $h(z)=\sum_KH(K)z^K$。因為 $z^N=z^{M_1}z^{N-M_1}$，$z$ 的次方會自己相加，所以
$\sum_Nz^N\sum_{M_1=0}^{N}\allowbreak H_1(M_1)H_2(N-M_1)=h_1(z)h_2(z)$。
\textbf{例}：$H_1=H_2=2^K=1,2,4,8,16$（一維 $K$ 步的走法數），摺積是 $1,4,12,32,80=(N+1)2^N$，
而 $h_1h_2=\frac1{(1-2z)^2}$；$z=0.1$ 時兩邊都是 1.5625（圖 \ref{seven-fig-dyson} 的表）。
Fourier 也一樣（(7.23)--(7.25)）：空間的摺積 $\int F_1(\vec x-\vec w)F_2(\vec w-\vec y)\,d^dw$ 變成乘積 $f_1(\vec q)f_2(\vec q)$。
\begin{errbox}[p.~176，(7.21)]
第二、三行印成 $Z^{N-M_1}$、$Z^{K_2}$（大寫），應為 $z^{N-M_1}$、$z^{K_2}$：大寫的 $Z$ 成了另一個變數，乘積就拆不開。
\end{errbox}

\textbf{一階的 Fourier 版}（(7.26)）：對步數取生成函數、對位置取 Fourier 之後，每段實線變成傳播子 $g_0(z;\vec q)=\frac1{1-z\chi(\vec q)}$（$\chi$ 是一步的結構函數），虛線變成 $F$：
\[g_0(z;\vec q)\Big[\sum_{\vec Q}F(\vec q-\vec Q)g_0(z;\vec Q)\Big]g_0(z;\vec q).\tag{7.26}\]
並排的兩個圈（Fig.~7.3(b)）就是兩個方括號 $g_0[\cdots]g_0[\cdots]g_0$（(7.27)）。
\textbf{費曼規則}（p.~180，摘要）：實線＝傳播子 $g_0$；虛線＝交互作用 $-F(\vec Q)$；每個頂點動量守恆；外動量 $\vec q$ 固定，其餘動量求和。
〔注〕規則說虛線是 $-F$，(7.26)--(7.34) 卻沒寫負號（把負號併進了 $F$）。
\begin{errbox}[p.~178，(7.27) 前]
書上說 (7.27) 是「Figure 7.3(a)」；(a) 是巢狀，(7.27) 是並排的 (b)。
\end{errbox}
\textbf{只用一階的圈、全部串起來}（(7.28)--(7.30)）：$\Sigma_1=\sum_{\vec Q}F(\vec q-\vec Q)g_0(z;\vec Q)$，
\[g\approx\sum_{n=0}^\infty g_0\big[\Sigma_1g_0\big]^n=\frac{1}{g_0^{-1}-\Sigma_1}=\frac1{1-z\chi(\vec q)-\Sigma_1(z;\vec q)}.\tag{7.30}\]

\subsubsection{Dyson 方程}
\leavevmode\bk{pp.~182--184}\ \anim{W704}\par\nopagebreak

\textbf{問}：所有的圖都要一張一張算嗎？能不能只算「剪不開的塊」，再自動串起來？

\textbf{自能}：剪斷一條實線就分成兩半的圖叫\textbf{可約}；剪不開的叫\textbf{不可約}（1PI）。不可約的圖（去掉兩端的傳播子）加起來就是自能 $\Sigma$；
$m$ 階有 $2m$ 個頂點、$m$ 條虛線、$2m-1$ 條實線、$m$ 個動量和。二階有兩種：巢狀 (7.34) 與交叉（Fig.~7.9、(7.31)），式子見練習 7.3。
〔注〕Fig.~7.10 畫的二階自能是交叉圖，(7.34) 寫的卻是巢狀圖；完整的二階自能兩個都要。
每一張圖都是一串不可約的塊用 $g_0$ 接起來，所以
\[g=g_0+g_0\Sigma g_0+g_0\Sigma g_0\Sigma g_0+\cdots=\frac{g_0}{1-\Sigma g_0}=\frac{1}{g_0^{-1}(z;\vec q)-\Sigma(z;\vec q)}.\tag{7.36, 7.37}\]
這是 \textbf{Dyson 方程}——像火車：車廂 $\Sigma$ 用連結器 $g_0$ 一節節接上（圖 \ref{seven-fig-dyson}）。
\textbf{算到底}：玩具數字 $g_0=2$、$\Sigma=\frac14$（所以 $\Sigma g_0=\frac12$）：部分和 $2,\ 3,\ \frac72,\ \frac{15}4,\ \frac{31}8\to4=\frac1{1/2-1/4}$。
\textbf{計數版}：把每條弧都換成數字 1，(7.37) 就變成「全部的弧線圖 $(2k-1)!!$＝不可約的 $1,2,10,74$ 排成一串」：
$\sum_{k\ge0}(2k-1)!!\,x^k=1/\big(1-(x+2x^2+10x^3+74x^4+\cdots)\big)$；右邊展開的係數是 $1,\ 1,\ 1+2=3,\ 3+2+10=15,\ 15+2\cdot3+10+74=105$。
\begin{errbox}[p.~183，(7.33)(7.35)]
(7.33) 第三項印成 $g_0[\Sigma_1+\Sigma_2]^2g_0$（少了中間的 $g_0$）；(7.35) 左邊印成 $\frac{g_0}{1+(\Sigma_1+\Sigma_2)g_0}$，應為 $1-$。
玩具數字照 (7.35) 得 $\frac43$、照 (7.33) 得 $\frac{10}3$，都不是 4。
\end{errbox}
〔注〕幾何級數只在 $|\Sigma g_0|<1$ 時收斂（一維接觸作用 $f=0.05$、$q=0.5$：$z=0.4$ 時 $|\Sigma g_0|=0.28$，收斂；$z=0.49$ 時是 1.80，發散），但閉式 (7.37) 仍然可用。

\begin{figure}[htbp]
\centering
\SevenFigDyson\hfill\SevenFigDysonBars\\[2pt]
{\footnotesize\setlength{\tabcolsep}{4pt}\SevenTabDyson}
\caption{生成函數與 Dyson 方程。左上：$g$ 是一串用實線 $g_0$ 接起來的自能 $\Sigma$；右上：玩具數字的部分和，逼近虛線 4。
左下：摺積變乘積的例子與 Dyson 的部分和（書上 (7.35) 的號給 $\frac43$）。右下：計數版的 Dyson 方程}\label{seven-fig-dyson}
\end{figure}

\subsubsection{微擾論為什麼失效}
\leavevmode\bk{pp.~184--186}\ \anim{W705}\par\nopagebreak

\textbf{問}：自能的每一階都帶一個小小的 $f$。是不是只要算前幾階就夠了？

\textbf{數冪次}（(7.38)(7.39)）：靠近 $z_c$、$q$ 小時 $g_0\approx\frac{1}{(z_c-z)+q^2}$（係數略去）。$l$ 階自能有 $2l-1$ 條傳播子、$l$ 個 $d^dq$ 積分，
$f^l\int\big(\frac1{(z_c-z)+q_i^2}\big)^{2l-1}\big(d^dq_i\big)^l$（(7.38)）。令 $q_i\to(z_c-z)^{1/2}q_i$：每條傳播子給 $(z_c-z)^{-1}$、每個積分給 $(z_c-z)^{d/2}$，
\[(z_c-z)^{-(2l-1)+dl/2}=(z_c-z)^{1-l(4-d)/2}.\tag{7.39}\]
圖 \ref{seven-fig-mf} 左上列出 $d=1..5$、$l=0..4$。$d=3$：$1,\frac12,0,-\frac12,-1$——每多一階冪次少 $\frac12$，高階比低階更奇異，$f$ 再小也壓不住；
$d=4$ 全是 1；$d=5$ 反而越來越不奇異。\textbf{四維是分界。}數值積分量出來的冪次：$l=1,d=1$：$-0.50$；$l=1,d=3$（截斷後的奇異部分）：$0.50$；
$l=2,d=1$（巢狀）：$-2.00$，和表一致。
\textbf{一維一階自能算到底}：$\int_{-\pi}^{\pi}\frac{dq}{2\pi}\frac1{1-2z\cos q}=\frac1{\sqrt{1-4z^2}}$；$z=0.4,\ 0.49,\ 0.499$ 得 1.6667、5.0252、15.8193：
離 $z_c=\frac12$ 每近 10 倍就大 $\sqrt{10}$ 倍（冪次 $-\frac12$）。
\textbf{後果}（(7.40)--(7.42)）：展開真正的參數是 $f(z_c-z)^{-(4-d)/2}$，$d<4$ 時 $z\to z_c$ 它趨於無窮；截到 $n_h$ 階，$g(z;0)$ 的奇點 $(z_c-z)^{-p}$、
$p=1-n_h(4-d)/2$ 每多算一階就變——低階微擾對長漫步一定失效。〔注〕(7.42) 第一行書上寫 $\Sigma(z_c-z;0)$，意思是 $\Sigma(z;0)$。

\subsubsection{平均場}
\leavevmode\bk{pp.~186--191}\ \anim{W705}\par\nopagebreak

\textbf{問}：既然一階一階加不起來，能不能找一個粗糙、但一次顧到全部的近似？

\textbf{模態展開}（(7.43)--(7.50) 的摘要）：在體積 $V$ 裡 $G(z;\vec x,\vec y)=\sum_i\psi_i(\vec x)\frac{z_c}{z_i-z}\psi_i^*(\vec y)$，模態是平面波 $\psi_{\vec q}=e^{i\vec q\cdot\vec x}/\sqrt V$，
$z_{\vec q}\approx z_c(1+z_ca^2q^2)$；連續極限滿足 $(1-z/z_c)\,G-za^2\nabla^2G=a^d\delta(\vec x-\vec y)$（(7.47)）。
\begin{errbox}[p.~188，(7.47)(7.48)、(7.50) 後]
(7.47) 印成 $za^2\nabla^2G-(z/z_c-1)G=a^d\delta$，$\nabla^2$ 那一項的號錯了：一維格子的精確 $G$ 代入，$z=0.4$、$x=1$ 兩邊差 0.3333，更正版差 0。
(7.48) 的 $\psi_i=e^{-\vec q\cdot\vec r}$ 應為 $e^{i\vec q\cdot\vec r}/\sqrt V$（否則 $\nabla^2\psi=+q^2\psi$，和 $\lambda_i=q^2$ 矛盾）。「equivalent to (7.43dd)」應為 (7.43d)。
\end{errbox}
\textbf{只留一個模態}：$z\to z_c$ 時最低的模態 $\psi_0=1/\sqrt{\mathcal N}$（$\mathcal N$ 個格點，$z_0=z_c$）主導，$G\approx G_0=z_c\frac{\psi_0(\vec x)\psi_0^*(\vec y)}{z_0-z}$（(7.51)）。每插入一條虛線是一個算子 $O=-\frac f2\int(\psi\psi^*)^2d^dw\,\frac{d^2}{dz^2}$（(7.52)），全部插入是 $e^{O}$（(7.54)）；把 $\frac{z_c}{z_0-z}=\int_0^\infty dt\,e^{-t(z_0-z)/z_c}$ (7.55) 代入，
$\frac{d^2}{dz^2}$ 就變成乘 $t^2$：
\[G_{MF}(z)=\frac1{\mathcal N}\int_0^\infty\exp\Big[-\frac{ft^2}{\mathcal N}-t(1-z/z_c)\Big]dt=\int_0^\infty\exp\big[\mathcal N\big(-f\lambda^2-\lambda(1-z/z_c)\big)\big]d\lambda .\tag{7.57}\]
〔注〕$\frac{d^2}{dz^2}e^{tz/z_c}=\frac{t^2}{z_c^2}e^{tz/z_c}$，所以 (7.52) 要多一個 $z_c^2$ 才得到 (7.53)(7.56)；(7.57) 又把 $\frac12$ 併進 $f$。$f$ 是唯象參數，結論不變。
\begin{errbox}[p.~189--190，(7.55) 前、(7.58)]
p.~189 說這個積分是 $1/(z_0-z)$；其實是 $z_c/(z_0-z)$（$z_c=\frac14$、$z_0=z_c$、$z=\frac15$：積分是 5，不是 20）。
p.~190 寫「展開 $e^{-\mathcal Nz\lambda/z_c}$」，(7.58) 第一行寫 $z^m$、$\lambda^n$；應為 $e^{+\mathcal Nz\lambda/z_c}$、$(z/z_c)^m$、$\lambda^m$。
\end{errbox}
\textbf{取出 $z^m$ 的係數}（(7.58)--(7.60)）：展開 $e^{+\mathcal N\lambda z/z_c}$，
\[C_{MF}(m)=\frac{1}{z_c^m}\frac{\mathcal N^m}{m!}\int_0^\infty e^{-\mathcal N(f\lambda^2+\lambda)+m\ln\lambda}d\lambda\approx\frac1{\mathcal N}z_c^{-m}e^{-fm^2/\mathcal N}.\tag{7.59, 7.60}\]
$f=0$ 時積分正好是 $m!/\mathcal N^{m+1}$，回到普通漫步的 $z_c^{-m}/\mathcal N$（$(2d)^m$ 條平均分到 $\mathcal N$ 個終點）；自迴避再乘 $e^{-fm^2/\mathcal N}$：
約 $\frac{m^2}2$ 對點、每對碰在一起的機率約 $\frac1{\mathcal N}$。\textbf{算到底}（練習 7.4，$fm^2/\mathcal N=1$）：$\mathcal Nz_c^mC_{MF}(m)$ 的精確積分在
$(m,\mathcal N)=(10,100)$：0.3385；$(100,10^4)$：0.3643；$(1000,10^6)$：0.3675，趨近 $e^{-1}=0.3679$（圖 \ref{seven-fig-mf} 右上）。
\begin{errbox}[p.~190，(7.59)]
指數印成 $-\mathcal N(f\lambda^2+\mathcal N\lambda)$，應為 $-\mathcal N(f\lambda^2+\lambda)$：照書上，$f=0$ 時多出 $\mathcal N^{-(m+1)}$ 倍，回不到普通漫步的 $z_c^{-m}/\mathcal N$。
\end{errbox}
\textbf{$z<z_c$ 時看不到自迴避}（練習 7.5）：$z/z_c=\frac12$ 時 $\mathcal N(1-z/z_c)G_{MF}$ 在 $\mathcal N=10,10^2,10^3,10^4$ 是 0.6899、0.9341、0.9922、0.9992，
趨近 1，和沒有自迴避一樣；$z/z_c=\frac32$ 時 $\ln G_{MF}/\mathcal N$ 是 0.0452、0.0596、0.0621，趨近 $\frac{(z/z_c-1)^2}{4f}=0.0625$。自迴避只在 $z\ge z_c$（長漫步主導）時才顯出來。

\subsubsection{Flory 的論證}
\leavevmode\bk{pp.~191--192}\ \anim{W706}\par\nopagebreak

\textbf{問}：$m$ 個單體的高分子（自迴避漫步），大小 $R$ 怎麼隨 $m$ 變大？普通漫步是 $R\propto m^{1/2}$。

\begin{figure}[htbp]
\centering
\begin{minipage}[c]{0.3\linewidth}\centering\small\setlength{\tabcolsep}{4pt}\SevenTabPower\end{minipage}\hfill
\begin{minipage}[c]{0.68\linewidth}\centering\footnotesize\setlength{\tabcolsep}{3.5pt}\SevenTabMF\end{minipage}\\[8pt]
\begin{minipage}[c]{0.42\linewidth}\centering\SevenFigFlory\end{minipage}\hfill
\begin{minipage}[c]{0.56\linewidth}\centering\footnotesize\setlength{\tabcolsep}{4pt}\SevenTabFlory\end{minipage}
\caption{微擾、平均場與 Flory。左上：$l$ 階自能的冪次 $1-l(4-d)/2$（(7.39)）。右上：練習 7.4、7.5 的精確積分。
左下：$d=2$、$m=1000$、$u=1$ 的熵（藍）、排斥（橙）與和（綠），最大在 $R^*$。右下：Flory 指數、數值最大化的斜率與已知值；二維精確列舉的局部斜率 $2\nu_{\rm eff}$}\label{seven-fig-mf}
\end{figure}

\textbf{排斥}（(7.61)）：把平均場用在半徑 $R$ 的一團上：$\mathcal N=\alpha R^d$、$u=f/\alpha$、$z_c=\frac1{2d}$，得
$C_{MF}(m;\vec x,\vec y)=\frac{(2d)^m}{\mathcal N}e^{-um^2/R^d}$——$R$ 越大越好。
\begin{errbox}[p.~191，(7.61)]
左邊印成 $C_{MF}(N;\vec x,\vec y)$，應為 $C_{MF}(m;\vec x,\vec y)$：步數是 $m$，而書上 §7.2.6 起的 $N$ 是格點數。
\end{errbox}
\textbf{熵}（(7.62)）：普通漫步走到 $R$ 的機率 $\propto e^{-R^2/m}$——$R$ 越小越好。兩者相乘就是拔河，最可能的 $R$ 讓指數最大：
\[\frac{d}{dR}\Big(-\frac{R^2}m-\frac{um^2}{R^d}\Big)=0\ \Rightarrow\ R^{d+2}=\frac{du}{2}m^3,\qquad R\propto m^{\nu_F},\quad\nu_F=\frac3{d+2}.\tag{7.63, 7.64}\]
\textbf{算到底}（$d=2$、$m=1000$、$u=1$，圖 \ref{seven-fig-mf} 下）：$R^{*4}=10^9$，$R^*=177.83$；此時兩項都是 $-31.62$、和是 $-63.25$
（$d=2$ 時極值條件 $\frac{2R}m=\frac{2um^2}{R^3}$ 正好就是「兩項相等」）。$d=1,2,3,4$：$\nu_F=1,\frac34,\frac35,\frac12$（$m=10^3\to10^6$ 數值最大化的斜率 1.0000、0.7500、0.6000、0.5000）。
已知值（只引用）：一維 1、二維 $\frac34$（精確）、三維約 0.588（數值）、四維以上 $\frac12$（四維有對數修正）。〔注〕熵因子正確是 $e^{-dR^2/2m}$（第 2 章的高斯），係數不影響指數。

\textbf{和精確計數比}：二維精確列舉的局部斜率 $2\nu_{\rm eff}=\ln(\langle R^2\rangle_N/\langle R^2\rangle_{N-2})/\ln\frac N{N-2}$，$N=6,8,10,12,14$：1.4306、1.4372、1.4442、1.4505、1.4556，
慢慢爬向 $2\nu=\frac32$（奇數 $N$：1.4459、1.4528、1.4562、1.4593）。\textbf{揭曉}：熵和排斥拔河，拔出 $R\propto m^{3/(d+2)}$；二維剛好是精確的 $\frac34$，三維的 0.6 也很接近 0.588。

% =====================================================================
\subsection{練習 7.1--7.5 解答}
{\small
\begin{enumerate}[label=\textbf{7.\arabic*},itemsep=0.05em]
\item 高斯的摺積。\quad 答：三個分量各自配方：$\int e^{-\alpha(x-r)^2-\beta(r-y)^2}dr=\sqrt{\frac{\pi}{\alpha+\beta}}\,e^{-\frac{\alpha\beta}{\alpha+\beta}(x-y)^2}$，
$\alpha=\frac{3}{2M_1a^2}$、$\beta=\frac3{2M_2a^2}$，所以 $\frac{\alpha\beta}{\alpha+\beta}=\frac{3}{2(M_1+M_2)a^2}$；前因子 $\sqrt{\frac\alpha\pi}\sqrt{\frac\beta\pi}\sqrt{\frac{\pi}{\alpha+\beta}}=\big(\frac{2\pi Na^2}{3}\big)^{-1/2}$，
$N=M_1+M_2$。三個分量相乘得 $\big(\frac{2\pi Na^2}{3}\big)^{-3/2}e^{-3|\vec x-\vec y|^2/2Na^2}$，對段數歸納得 $k$ 段；前因子必須用 $(2\pi Ma^2/3)^{-3/2}$（書上少 3）。
\item 三階的圖。\quad 答：6 個點配成 3 條弧，$5\cdot3\cdot1=15$ 種（圖 \ref{seven-fig-chord}）。不可約 10 種；可約 5 種：三個一階的圈並排 1 種（$1\cdot1\cdot1$），
一階圈接二階不可約圖（巢狀或交叉）2 種（$1\cdot2$），倒過來 2 種（$2\cdot1$）。若把頭尾倒過來的算同一種，剩 11 種（二階仍是 3 種）。
\item 二階的三張圖。\quad 答（每條虛線 $-F$，兩條相乘為正）：(a) 巢狀 $g_0(\vec q)\Big[\sum_{\vec Q_1,\vec Q_2}F(\vec q-\vec Q_1)g_0(\vec Q_1)F(\vec Q_1-\vec Q_2)g_0(\vec Q_2)g_0(\vec Q_1)\Big]g_0(\vec q)$；
(b) 並排 $g_0\Big[\sum_{\vec Q}F(\vec q-\vec Q)g_0(\vec Q)\Big]g_0\Big[\sum_{\vec Q}F(\vec q-\vec Q)g_0(\vec Q)\Big]g_0$（可約，不屬於自能）；
(c) 交叉 $g_0(\vec q)\Big[\sum_{\vec Q_1,\vec Q_2}F(\vec q-\vec Q_1)F(\vec Q_1-\vec Q_2)g_0(\vec Q_1)g_0(\vec q+\vec Q_2-\vec Q_1)g_0(\vec Q_2)\Big]g_0(\vec q)$。
\item 平均場的漸近式。\quad 答：令 $t=\mathcal N\lambda$，$\mathcal Nz_c^mC_{MF}(m)=\frac1{m!}\int_0^\infty t^me^{-t}e^{-ft^2/\mathcal N}dt=\langle e^{-fT^2/\mathcal N}\rangle$，
$T$ 是 Gamma$(m+1)$ 分布（平均 $m+1$、標準差 $\sqrt{m+1}$）。$m$ 大時 $T\approx m$，得 $e^{-fm^2/\mathcal N}$，即 (7.60)；修正約 $f^2m^3/\mathcal N^2$
（$fm^2/\mathcal N=1$ 時相對修正約 $\frac1m$：表裡的 0.3385、0.3643、0.3675 比 $e^{-1}$ 小 8.0\%、1.0\%、0.1\%）。
\item $z<z_c$ 與 $z>z_c$。\quad 答：記 $\varepsilon=1-z/z_c>0$，指數在 $\lambda=0$ 最大，展開 $e^{-\mathcal Nf\lambda^2}$ 得 $G_{MF}\approx\frac{1}{\mathcal N\varepsilon}-\frac{2f}{\mathcal N^2\varepsilon^3}$；
主項 $\frac1{\mathcal N}\sum_m(z/z_c)^m$ 表示條數是 $z_c^{-m}/\mathcal N$，自迴避沒有效果；$\mathcal N\varepsilon G_{MF}\approx1-\frac{2f}{\mathcal N\varepsilon^2}$ 在 $\mathcal N=10^3,10^4$ 是 0.992、0.9992，和表一致。
$z>z_c$ 時最大在 $\lambda^*=\frac{z/z_c-1}{2f}$，$G_{MF}\sim e^{\mathcal N(z/z_c-1)^2/4f}$，所以 $\ln G_{MF}/\mathcal N\to0.0625$。
\end{enumerate}}

% =====================================================================
\subsection{第 7 章的回顧與筆誤一覽}
\begin{center}
\small\setlength{\tabcolsep}{4pt}
\begin{tabular}{@{}llll@{}}
\toprule
主題 & 結果 & 書上的式子 & 動畫\\
\midrule
路徑和 & 切段求和：4 步回原點一維 6、二維 36；高斯變異數相加 & (7.1)--(7.8) & W701\\
自迴避 & $c_{14}=2374444$、存活 0.0088、$\langle R^2\rangle=42.79$ 對普通漫步的 14 & (7.9)(7.10) & W702\\
$1-f$ 展開 & 弧線圖 $(2k-1)!!$；$N=4$ 部分和 256、28、124、96、100 & (7.11)--(7.19) & W703\\
Dyson & 摺積變乘積；$g=1/(g_0^{-1}-\Sigma)$；計數版 $1,2,10,74\to1,3,15,105$ & (7.20)--(7.37) & W704\\
微擾、平均場 & 冪次 $1-l(4-d)/2$，四維是分界；$C_{MF}\approx\frac1{\mathcal N}z_c^{-m}e^{-fm^2/\mathcal N}$ & (7.38)--(7.61) & W705\\
Flory & $R\propto m^{3/(d+2)}$：二維 $\frac34$（精確）、三維 0.6（數值 0.588） & (7.62)--(7.64) & W706\\
\bottomrule
\end{tabular}
\end{center}
\textbf{筆誤一覽}（給 2004 年版；〔注〕見各節）：下表在文中都用紅框標出，每一條都用程式核對過。
\begin{center}
\footnotesize\renewcommand{\arraystretch}{0.9}
\SevenTabErrata
\end{center}
